%=============================================================================!
% 4.5  LARGER SWINGS: ANHARMONIC WELLS                                        !
%=============================================================================!
\section{Larger swings: anharmonic wells}
\label{sec:os-anharmonic}

Real wells are not parabolas. The pendulum's well rises ever more slowly
than its parabola as the bob swings towards the horizontal, and the well that holds two atoms together
rises steeply when they are pushed together but levels off when they are
pulled apart. A well that is not a parabola is called anharmonic. This
section computes the period of a swing of any size from the energy alone,
and finds that in every well of interest here it grows with the energy.

\subsection*{The period from the energy}

\Cref{sec:en-track} noted that energy fixes the time of a journey through
the speed: a short step $\dd x$ taken at speed $v$ lasts $\dd x/v$. In a
well, with total energy $E$ and turning points $x_1 < x_2$, the speed at
$x$ is $v(x) = \sqrt{2(E - U(x))/m}$, so the journey from $x_1$ to $x_2$
takes $\int_{x_1}^{x_2}\dd x/v(x)$. The way back takes as long, since the
speed at each point is the same in either direction, and the period is
\begin{equation}
   \mathcal{T}(E) = 2\int_{x_1}^{x_2}\frac{\dd x}{\sqrt{2\bigl(E - U(x)\bigr)/m}} .
   \label{eq:os-period}
\end{equation}
At the turning points the speed is zero and the integrand is infinite.
The integral then means the limit as its ends approach the turning
points, and it is finite when, as in the harmonic case below, the
function whose derivative is the integrand tends to finite values there. To evaluate it we need
the derivative of the arcsine.

\begin{toolbox}[The derivative of an inverse function]
For $u$ between $-1$ and 1, $\theta = \arcsin u$ is the angle between
$-\pi/2$ and $\pi/2$ whose sine is $u$ (\cref{sec:mo-trig}), so $\sin
\theta(u) = u$ for every such $u$. Differentiating both sides with respect
to $u$, the left by the chain rule,
\begin{equation*}
   \cos\theta\,\frac{\dd\theta}{\dd u} = 1 ,
   \qquad\text{so}\qquad
   \frac{\dd}{\dd u}\arcsin u = \frac{1}{\cos\theta}
   = \frac{1}{\sqrt{1 - u^2}} ,
\end{equation*}
since $\cos\theta = +\sqrt{1 - \sin^2\theta}$ for an angle between
$-\pi/2$ and $\pi/2$, where the cosine is not negative. With $u = x/A$ and
the chain rule, $\dd u/\dd x = 1/A$, so
\begin{equation*}
   \frac{\dd}{\dd x}\arcsin\frac{x}{A} = \frac{1}{A}\,
   \frac{1}{\sqrt{1 - x^2/A^2}} = \frac{1}{\sqrt{A^2\,(1 - x^2/A^2)}}
   = \frac{1}{\sqrt{A^2 - x^2}} ,
\end{equation*}
the positive $A$ having gone inside the root as $A^2$.
The same method, differentiating the equation that defines the inverse,
gives the derivative of any inverse function.
\end{toolbox}

For the spring, $U = \frac12 kx^2$ and $E = \frac12 kA^2$, so $E - U =
\frac12 k(A^2 - x^2)$ and $v = \omega\sqrt{A^2 - x^2}$. Then
\begin{align*}
   \mathcal{T} &= \frac{2}{\omega}\int_{-A}^{A}\frac{\dd x}{\sqrt{A^2 - x^2}}
   = \frac{2}{\omega}\Bigl[\arcsin\frac{x}{A}\Bigr]_{-A}^{A}
      && \text{by the toolbox and \cref{eq:mo-integrals}}\\
   &= \frac{2}{\omega}\Bigl(\frac{\pi}{2} - \Bigl(-\frac{\pi}{2}\Bigr)\Bigr)
   = \frac{2\pi}{\omega} ,
\end{align*}
the period of \cref{sec:os-circle}, whatever the amplitude. The integrand
is infinite at $\pm A$, yet $\arcsin(x/A)$ rises only to $\pm\pi/2$ there,
so the area under it is finite: the spike is tall but very narrow. For
other wells the integral is evaluated by computer. The function
\texttt{oscillators.period} of \texttt{mdlab} finds the turning points,
then writes $x = c - \eta\cos\vartheta$, with $c$ the centre and $\eta$ the
half-width of the allowed region, so that $\vartheta$ runs from 0 at
$x_1$ to $\pi$ at $x_2$. A short step $\dd\vartheta$ moves $x$ by $\dd x =
\eta\sin\vartheta\,\dd\vartheta$, the derivative of $c - \eta\cos\vartheta$
times the step, so each time $\dd x/v$ becomes $\eta\sin\vartheta\,
\dd\vartheta/v$, and these are added up over $\vartheta$ from 0 to $\pi$.
Near $x_1$, where $U(x_1) = E$, the Taylor series gives $E - U(x) \approx
-U'(x_1)\,(x - x_1)$, and $x - x_1 = \eta(1 - \cos\vartheta) \approx \frac12
\eta\vartheta^2$, so $v$ grows in proportion to $\vartheta$; so does
$\sin\vartheta \approx \vartheta$, and their ratio, the new integrand,
stays finite, as it does near $x_2$. The integral over $\vartheta$ is
then evaluated by Gauss--Legendre quadrature, a weighted sum of the
integrand at chosen points, which converges quickly for a smooth
integrand. Its tests check it against the spring, against closed
forms for the pendulum and for the Morse well below, and against motions
timed by \texttt{solve\_newton}.

\subsection*{The pendulum and the hop}

The pendulum's well, $U = mgL(1 - \cos\theta)$, and the hop coordinate of
the lithium atom, \cref{eq:os-hop}, have the same shape,
\begin{equation*}
   U = \tfrac12 D\,\Bigl(1 - \cos\frac{2\pi x}{\Lambda}\Bigr) ,
\end{equation*}
a cosine well of depth $D$ that repeats every $\Lambda$: for the
pendulum $x$ is the distance $s = L\theta$ along the arc, $\Lambda =
2\pi L$ and $D = 2mgL$, the potential energy of the bob held upright; for the atom $\Lambda = a$ and $D = U_{\mathrm b}$,
the barrier. Away from its bottom the cosine well rises more slowly than
its parabola (\cref{fig:os-wells}a), so the restoring force there is
weaker than a spring's, and in a swing of the same amplitude the body
moves more slowly and takes longer. \Cref{eq:os-period} gives, for the pendulum,
\begin{center}
\begin{tabular}{lccccc}
\toprule
amplitude & \ang{10} & \ang{45} & \ang{90} & \ang{150} & \ang{179}\\
\midrule
period / small-swing period & 1.0019 & 1.0400 & 1.1803 & 1.7622 & 3.9011\\
\bottomrule
\end{tabular}
\end{center}
For swings of \ang{10} the period is within \SI{0.2}{\percent} of
$2\pi\sqrt{L/g}$, and it changes only slowly with the amplitude: a clock
whose swing grows from \ang{10} to \ang{11} loses about \SI{35}{\second} a
day, which is why clockmakers keep the swing small and steady. Beyond \ang{90} the bob
ends each swing above the pivot, where only a rod, which can push, keeps
it on its circle; a string would go slack. As the amplitude approaches
\ang{180}
the bob lingers ever longer near the top, where the force on it is
nearly zero, and the period grows without limit (\cref{fig:os-period}).

The lithium atom behaves the same way. With \SI{0.25}{\electronvolt} above
the bottom of the hollow,
below the \SI{0.3}{\electronvolt} barrier, it rattles in its hollow with a
period of \SI{253.8}{\femto\second}, half as long again as the
\SI{170.3}{\femto\second} of small swings, and with
\SI{0.297}{\electronvolt} the period is \SI{400.8}{\femto\second}. With exactly
the barrier as energy, the atom would creep ever more slowly towards the
bridge without ever reaching it. The phase
portrait of the cosine well (\cref{fig:os-phase}b) shows all of this at
once: closed paths for swings or rattles, the path through the top that
separates them from the open paths of a pendulum going over the top or an
atom hopping on from hollow to hollow.

\subsection*{Two wells for a pair of atoms}

The energy of two atoms depends on the distance $r$ between them, and two
standard shapes model it; Chapter~8 discusses where they apply. The Morse
well, used for a chemical bond, is
\begin{equation*}
   U_{\mathrm{M}}(r) = D\,\bigl(1 - \mathrm{e}^{-a(r - r_0)}\bigr)^2 ,
\end{equation*}
with three constants. At $r = r_0$ the exponential is 1 and $U_{\mathrm M}
= 0$, its least value, since it is a square; when the atoms are pulled
far apart the exponential vanishes and $U_{\mathrm M}$ levels off at $D$,
the energy needed to separate them; when they are pushed together,
$r < r_0$, the exponential grows rapidly and so does $U_{\mathrm M}$. The
constant $a$, an inverse length, sets the width of the well; here it is
not the spacing of the surface. Near the bottom, with $y = r - r_0$ and
the exponential series of \cref{sec:ne-drag}, $\mathrm{e}^{-ay} = 1 - ay +
\order{y^2}$, so $1 - \mathrm{e}^{-ay} = ay + \order{y^2}$ and, squaring,
\begin{equation*}
   U_{\mathrm M} = D a^2 y^2 + \order{y^3} ,
   \qquad\text{a spring of stiffness } k = 2Da^2 .
\end{equation*}

The Lennard-Jones well, often used for atoms that form no bond, such as those
of the noble gases (helium, neon, argon and their family, which hardly
react), is
\begin{equation*}
   U_{\mathrm{LJ}}(r) = 4\varepsilon\Bigl[\Bigl(\frac{\sigma}{r}\Bigr)^{12}
   - \Bigl(\frac{\sigma}{r}\Bigr)^6\Bigr] ,
\end{equation*}
a steep repulsion, $r^{-12}$, that dominates at short range and a weaker
attraction, $-r^{-6}$, that dominates at long range. The power rule of
\cref{sec:mo-velocity} was found for whole $n \ge 1$, but it holds for
$r^{-n} = 1/r^n$ too: differentiating $r^{-n}\,r^n = 1$ by the product
rule gives $r^n\,\dd(r^{-n})/\dd r + r^{-n}\,n r^{n-1} = 0$, so $\dd
(r^{-n})/\dd r = -n\,r^{-n-1}$. Differentiating with this rule,
\begin{equation*}
   \frac{\dd U_{\mathrm{LJ}}}{\dd r}
   = 4\varepsilon\Bigl[-12\,\frac{\sigma^{12}}{r^{13}}
   + 6\,\frac{\sigma^6}{r^7}\Bigr]
   = \frac{24\varepsilon}{r}\Bigl[\Bigl(\frac{\sigma}{r}\Bigr)^6
   - 2\Bigl(\frac{\sigma}{r}\Bigr)^{12}\Bigr] ,
\end{equation*}
which vanishes when the bracket, $(\sigma/r)^6\,[1 - 2(\sigma/r)^6]$,
does, so when $(\sigma/r)^6 = \frac12$, that is $(r/\sigma)^6 = 2$, at
$r_{\min} = 2^{1/6}\sigma = 1.1225\,\sigma$, where $U_{\mathrm{LJ}} = 4\varepsilon(\frac14 - \frac12)
= -\varepsilon$. Differentiating once more,
\begin{equation*}
   \frac{\dd^2 U_{\mathrm{LJ}}}{\dd r^2}
   = \frac{4\varepsilon}{r^2}\Bigl[156\Bigl(\frac{\sigma}{r}\Bigr)^{12}
   - 42\Bigl(\frac{\sigma}{r}\Bigr)^6\Bigr]
   = \frac{4\varepsilon}{r_{\min}^2}\Bigl[\frac{156}{4} - \frac{42}{2}\Bigr]
   = \frac{4\times18\,\varepsilon}{2^{1/3}\sigma^2}
   = 57.15\,\frac{\varepsilon}{\sigma^2}
\end{equation*}
at the minimum, the middle steps putting in $(\sigma/r)^{12} = \frac14$,
$(\sigma/r)^6 = \frac12$ and $r_{\min}^2 = 2^{2/6}\sigma^2 =
2^{1/3}\sigma^2$, with $39 - 21 = 18$.

\begin{figure}[htb]
\centering
\includegraphics{ch04_oscillations/wells}
\caption{Three anharmonic wells (black) against the parabolas of the same
curvature at their bottoms (dashed), each in its own natural units; the
dotted lines mark half the depth. (a) The cosine well of the pendulum and
of the lithium hop. (b) The Morse well of a bond. (c) The Lennard-Jones
well, drawn from its bottom.}
\label{fig:os-wells}
\end{figure}

Both wells are lopsided (\cref{fig:os-wells}b, c): steeper than their
parabola on the inside, softer on the outside. A swing therefore reaches
less far in and much further out than a spring's would; in the
Lennard-Jones well with half the depth as energy, the turning points are
at $1.027\,\sigma$ and $1.377\,\sigma$, against $0.990\,\sigma$ and
$1.255\,\sigma$ for the parabola, whose turning points lie
$\sqrt{2\times0.5/57.15} = 0.132\,\sigma$ either side of the minimum at
$1.1225\,\sigma$. The soft outer side wins, and the period grows with the
energy (\cref{fig:os-period}). Measured from the bottom, at a tenth, half
and nine tenths of the depth the period is 1.054, 1.414 and 3.162 times
that of small swings for the Morse well, which are $1/\sqrt{1 - E/D}$
exactly, the closed form the tests use, and 1.069, 1.554 and 4.399 times
for the Lennard-Jones well. These ratios hold for any constants and any
mass, because in the units of the figures the wells have no constants
left, and by \cref{eq:os-period} every period is proportional to
$\sqrt m$. A bond with
more energy therefore vibrates more slowly, and the period of the small
swings, the one that sets the steps of a simulation, is the shortest.

\begin{figure}[htb]
\centering
\includegraphics{ch04_oscillations/period}
\caption{The period of the motion in each well of \cref{fig:os-wells},
divided by the period of small swings, against the energy above the
bottom as a fraction of the depth, computed from \cref{eq:os-period}. The
cosine well's period grows without limit as the energy approaches the
barrier, though so slowly that it is only 3.8 at 0.9999 of the barrier,
where its curve stops; the others' grows without limit as the energy
approaches that needed to separate the atoms. A spring keeps the ratio 1 (dotted).}
\label{fig:os-period}
\end{figure}

\notebook{04}{explore the period in each well as the energy grows}

\begin{selfcheck}
Why does a pendulum swinging out to \ang{90} on each side take longer for
each swing than one swinging out to \ang{10}, even though it moves faster at the
bottom?
\end{selfcheck}
